% RnE 양식 구현을 위한 프리엠블 2024.08.21.

\usepackage{fapapersize}
\usefastocksize{210mm,297mm}
\usefapapersize{210mm,297mm,25mm,25mm,25mm,25mm} 
\usepackage{mathtools}
\usepackage{amsthm}
\usepackage{amssymb}
\usepackage{tikz}
\usetikzlibrary{calc,decorations.pathmorphing,shapes,patterns}
\usepackage{rotating}
\usepackage{pgf,pgfplots}
\pgfplotsset{compat=1.15}
\usepackage{mathrsfs}
\usetikzlibrary{arrows}
\usepackage{lipsum}
\usepackage{jiwonlipsum}
\usepackage{tabularray}
\UseTblrLibrary{varwidth,booktabs}
\usepackage{ragged2e}   
\usepackage{textpos}
\usepackage{kskorsymb}

\renewcommand\thesection{\Roman{section}}
\renewcommand\thesubsection{\arabic{subsection}}
\renewcommand\thesubsubsection{\thesubsection.\arabic{subsubsection}}
\setsecnumformat{\csname the#1\endcsname.\ }

%\linespread{1.3} % 영어 논문이므로 줄간격 조절함

\usepackage{caption}

\DeclareCaptionLabelFormat{cjk-single-bracket}{[#1 #2]}
\DeclareCaptionLabelFormat{cjk-angle-bracket}{<#1 #2>}
\captionsetup[figure]{%
labelsep=space,
labelfont=bf,
labelformat=cjk-single-bracket,
name={Fig.}
}
\captionsetup[table]{%
labelsep=period,
labelfont=bf,
%labelformat=cjk-angle-bracket,
name={Table}
}


%%%%%%%%%%%%%%%%%%%%%%%%%%%%%%%%%%%%%%%%%%%%%%%%%%%%%%%%
%-- 정리류 환경 설정
\newtheorem{theorem}{\bfseries\sffamily 정리}%[section]
\newtheorem{lemma}[theorem]{도움정리}
\newtheorem{corollary}[theorem]{따름정리}
\theoremstyle{definition}
\newtheorem{definition}[theorem]{정의}
\newtheorem{example}[theorem]{보기}
\theoremstyle{remark}
\newtheorem{remark}[theorem]{참고}

\numberwithin{equation}{section}
\renewcommand*{\proofname}{증명}
%==
% \renewcommand\qedsymbol{
% \hfill\ensuremath{\blacksquare}
%} 
%====
%-- 정리류 환경 설정 끝

%---- description label 변경
% \renewcommand*\descriptionlabel
% [1]{\hspace\labelsep
% \scshape\mdseries #1}
%---------- description label 변경 끝

\newcommand{\splitatcommas}[1]{%
  \begingroup
  \begingroup\lccode`~=`, \lowercase{\endgroup
    \edef~{\mathchar\the\mathcode`, \penalty0 \noexpand\hspace{0pt plus 1em}}%
  }\mathcode`,="8000 #1%
  \endgroup
}

\newcommand{\mycheckedbox}{$\square $\!\!\!\!\checkmark}




\setcounter{secnumdepth}{3}
\renewcommand{\abstractname}{Abstract} %Abstract 이름 변경

\input{Ltrominoes}
\input{biblatexcontrol.tex} % 참고문헌 표기 방식 설정 파일

\allowdisplaybreaks

%%% RnE 보고서 관련 요약문 등 %%%

%%% 연구결과 요약문 및 과제 수행 과정의 표에 들어가는 내용을 쓰는 곳이다.

%%%--------------------------------------
\newcommand{\summaryGoal}{%%% 연구 필요성 및 목적 쓰는 곳
\jiwon[1] %% todo: 여기의 \jiwon[1]를 삭제하고 "연구의 필요성 및 목적"을 쓸 것 
% "연구 필요성 및 목적 쓰는 곳 끝.
}
%%%---------------------------------------


\newcommand{\summaryResult}{%%% 연구내용 및 연구결과 쓰는 곳
\jiwon[2-3] %% todo: 여기의 \jiwon[2-3]을 삭제하고 "연구내용 및 연구결과"를 쓸 것 
% 연구내용 및 연구결과 쓰는 곳 끝.
}

%%%---------------------------------------
\newcommand{\summaryEffect}{%%% 연구 기대효과 쓰는 곳
\jiwon[4] %% todo: 여기의 \jiwon[4]를 삭제하고 "연구 기대효과"를 쓸 것 
% 연구 기대효과 쓰는 곳 끝.
}

%%%---------------------------------------


\newcommand{\summaryExpert}{%%% 전문가 피드백 활용 과정 쓰는 곳
% 전문가 피드백 활용을 받지 않았다면, 아래 줄의 \jiwon[5]를 삭제하고 "해 당 없 음"이라고 쓴다.
\jiwon[5] %% todo: 전문가 피드백 활용과정에 쓸 내용이 있으면, \jiwon[5]를 지우고 내용을 쓰면 된다.
% 전문가 피드백 활용 과정 쓰는 곳 끝.
}

%%%---------------------------------------


\newcommand{\summaryProblemSolving}{%%% 문제해결 과정 쓰는 곳
\jiwon[6-7] %% todo: 여기의 \jiwon[6-7]을 삭제하고 "문제해결 과정"을 쓸 것 
% 문제해결 과정 쓰는 곳 끝.
}

%%%---------------------------------------
\newcommand{\summaryFirstMember}{%%% 팀원1 내용 쓰는 곳
\jiwon[8] %% todo: 여기의 \jiwon[8]을 삭제하고 "팀원1 "의 내용을 쓸 것 
% 팀원1 내용 쓰는 곳 끝.
}

%%%---------------------------------------

% \newcommand{\summarySecondMember}{%%% 팀원2 내용 쓰는 곳


% % 팀원2 내용 쓰는 곳 끝.
% }


%%%---------------------------------------
\newcommand{\summaryLearning}{%%% 과제 수행 과정에서의 학습 쓰는 곳
\jiwon[9] %% todo: 여기의 \jiwon[9]를 삭제하고 "과제 수행 과정에서의 학습"의 내용을 쓸 것 
% 과제 수행 과정에서의 학습 쓰는 곳
}



\begin{document}
\pagestyle{empty}

\begin{textblock*}{\textwidth}(0cm,0cm)
  \noindent\LARGE\bfseries 1. 연구결과 요약문
\end{textblock*}
\vspace{23pt}
\begin{center}
  \begin{tblr}{
      width = \textwidth,
      colspec = {Q[1,c,m,gray!20]Q[3.88,m]},
      hlines,
      vline{1-Z} = {solid},
      column{1}={font=\bfseries\sffamily}
      }
      과 제 명 & \mytitleko  \\
      {연구 필요성\\ 및 목적} & \summaryGoal \\
      {연구내용 \\ 및 \\ 연구결과} & \summaryResult  \\
      연구 기대효과& \summaryEffect  \\
      {Keyword\\ (5개 이내)}& \mykeywordsko \\
  \end{tblr}    
\end{center}

\newpage

\begin{textblock*}{\textwidth}(0cm,0cm)
  \noindent\LARGE\bfseries 2. 과제 수행 과정
\end{textblock*}
\vspace{23pt}
\begin{center}
  \newlength\mycellwidth
% \setlength\mycellwidth{.76\textwidth}
\setlength\mycellwidth{\dimexpr\linewidth-\parindent-2.5cm-2.5em\relax} % 셀 병합에서 사용한 길이 지정
  \begin{tblr}{
      width=\textwidth,
      measure=vbox,
      colspec = {Q[2.5cm,c,m,gray!20]X[.7,m]X[4.96,m]},
      hlines,
      vline{1-Z} = {solid},
      % cell{1}{2}={c=2}{},
      % cell{2}{1}={r=2}{},
      % cell{2}{2}={c=2}{},
      % cell{Z}{2}={c=2}{},
      column{1}={font=\bfseries\sffamily}
      }
      {전문가 \\ 피드백 활용 \\ 과정} & \SetCell[c=2]{c} {\hsize=\mycellwidth\justifying\summaryExpert\par} &  \\
      \SetCell[r=2]{c,m} {문제해결 \\ 과정} & \SetCell[c=2]{m}  {\hsize=\mycellwidth\justifying\summaryProblemSolving\par}  &  \\
      & \SetCell{l,t} 팀원1 & \summaryFirstMember \\
      % & & \\
      % & & \\
      {과제 수행 \\ 과정에서의 \\ 학습}& \SetCell[c=2]{c} {\hsize=\mycellwidth\justifying\summaryLearning\par}  &   \\
  \end{tblr}        
\end{center}

\newpage


\begin{textblock*}{\textwidth}(0cm,0cm)
  \noindent\LARGE\bfseries 3. 연구결과
  \end{textblock*}
  \vspace{18pt} 
\begin{center}
\begin{spacing}{1.3}
{\LARGE\bfseries    \mytitleko }\par
\vspace{5pt}
\textcolor{gray}{\myauthorfirstko} \par 
\textcolor{gray}{\affiliationko}
\end{spacing}
\vspace{10pt}

{\large\bfseries    \mytitle }\par
\vspace{5pt}
\textcolor{gray}{\myauthorfirst} \par 
\textcolor{gray}{\affiliation}
\end{center}